% !TeX root = ../QuestionSet.tex
% ==================== 第1章 集合与常用逻辑用语 ====================
\QuestionSetChapter{集合与常用逻辑用语}

% ===== 1.1 集合的概念与运算 =====
\QuestionSetSection{集合的概念与运算}

\questionGroup{一、选择题}

% ----------------------------------------
\begin{question}
设全集$U=\left\{ x\mid x<9 \right\}$，集合$A=\{1,3,5\}$，则$\complement_U A$中元素个数为\choiceBox
\choices{0}{3}{5}{8}
\end{question}
\begin{answer}{C}
全集 $U=\{1,2,3,4,5,6,7,8\}$，$\complement_UA=\{2,4,6,7,8\}$，共 $5$ 个元素。选 C。
\end{answer}

% ----------------------------------------
\begin{question}[GPT生成]
设集合$A=\{x\in\mathbf{Z}\mid -2\le x<3\}$，$B=\{-1,0,2,4\}$，则$A\cap B=$\choiceBox
\choices{$\{-1,0,2\}$}{$\{-2,-1,0,1,2\}$}{$\{4\}$}{$\{-1,0,1,2,4\}$}
\end{question}
\begin{answer}{A}
$A=\{-2,-1,0,1,2\}$，与 $B=\{-1,0,2,4\}$ 的公共元素为 $-1,0,2$，故 $A\cap B=\{-1,0,2\}$。选 A。
\end{answer}

% ----------------------------------------
\begin{question}[GPT生成]
已知集合$A=\{1,2,3\}$，$B=\{2,3,4,5\}$，则$A\cup B$中元素个数为\blank{3cm}．
\end{question}
\begin{answer}{5}
$A\cup B=\{1,2,3,4,5\}$，共有 $5$ 个元素。
\end{answer}

% ===== 1.2 常用逻辑用语 =====
\QuestionSetSection{常用逻辑用语}

\questionGroup{一、选择题}

% ----------------------------------------
\begin{question}[GPT生成]
设命题$p:x>2$，命题$q:x^2>4$，则$p$是$q$的\choiceBox
\choices[2]{充分不必要条件}{必要不充分条件}{充要条件}{既不充分也不必要条件}
\end{question}
\begin{answer}{A}
若 $x>2$，则 $x^2>4$，所以 $p\Rightarrow q$；但 $x^2>4$ 时可能 $x<-2$，不一定有 $x>2$，所以 $q\nRightarrow p$。故 $p$ 是 $q$ 的充分不必要条件。选 A。
\end{answer}

% ----------------------------------------
\begin{question}[GPT生成]
命题“存在实数$x$，使得$x^2+1=0$”的否定为\blank{4cm}．
\end{question}
\begin{answer}{对任意实数$x$，都有$x^2+1\ne0$}
存在量词命题的否定要把“存在”改为“任意”，并否定结论，因此否定为“对任意实数 $x$，都有 $x^2+1\ne0$”。
\end{answer}

\QuestionSetSection*{本章综合训练}

\questionGroup{综合解答题}

% ----------------------------------------
\begin{question}[GPT生成]
已知集合$A=\{x\mid x^2-3x+2=0\}$，$B=\{x\mid x^2-(m+1)x+m=0\}$。（1）求集合$A$；（2）若$B\subseteq A$，求实数$m$的取值；（3）若$A\cap B=\{1\}$，求$m$的取值范围。
\end{question}
\solutionSpace{5cm}
\begin{answer}{(1) $A=\{1,2\}$；(2) $m=1$或$m=2$；(3) $m\ne2$}
(1) 由 $x^2-3x+2=(x-1)(x-2)=0$，得 $A=\{1,2\}$。(2) $x^2-(m+1)x+m=(x-1)(x-m)$，故 $B=\{1,m\}$（当 $m=1$ 时为单元素集）。要使 $B\subseteq A$，需 $m\in\{1,2\}$。(3) 因 $1\in A\cap B$ 恒成立，要使交集只有 $1$，只需 $m\ne2$。
\end{answer}

% ----------------------------------------
\begin{question}[GPT生成]
设命题$p:-1\le x\le3$，命题$q:a-1\le x\le a+2$。（1）若$p$对应集合为$P$，$q$对应集合为$Q$，分别写出$P,Q$；（2）若$q$是$p$的充分条件，求$a$的取值范围；（3）当$a=2$时，判断$p$与$q$的关系。
\end{question}
\solutionSpace{5cm}
\begin{answer}{(1) $P=[-1,3]$，$Q=[a-1,a+2]$；(2) $0\le a\le1$；(3) 既不充分也不必要}
(1) 由题意 $P=[-1,3]$，$Q=[a-1,a+2]$。(2) $q$ 是 $p$ 的充分条件等价于 $Q\subseteq P$，故 $a-1\ge-1$ 且 $a+2\le3$，解得 $0\le a\le1$。(3) 当 $a=2$ 时，$Q=[1,4]$，与 $P=[-1,3]$ 互不包含，所以二者既不充分也不必要。
\end{answer}

\printChapterAnswers
